\section{Three exclusions by Galois-unit pairing and parity}
\label{sec:galois-pairing-20261027}
\subsection{The projective-parity exclusion of $R9\_2\_1975$}
We use the categorical TY/Morita reduction and the general-complex
sphericalization and Galois-existence gate established above. The extreme
Morita branch has already been excluded. In either uniform branch the
rank-twelve dual has invertibles $G=D_8$, central subgroup $H=\langle z\rangle
\cong C_2$, quotient $Q=G/H\cong V_4$, and four other simples $X_q$ with
\[
 gX_q=X_qg=X_{q+gH},\qquad
 X_pX_q=\sum_{gH=p+q+\alpha}g+2\sum_{t\in Q}X_t,
 \qquad X_q^*=X_{q+\alpha}.
\]
Their common dimension is $d=4+3\sqrt2$, with $d^2=2+8d$ and
$D=8+4d^2=144+96\sqrt2$. The shift $\alpha$ is a reflection or rotation
coset. The forced pointed associator is the TY-origin class described in
the preceding section. Its full class, not only its cyclic restrictions,
is fixed by the detection argument cited there.

\subsubsection{Duality and an arbitrary complex stabilizer cocycle}
Let $K=G\times G$ act by $(g,k)X=gXk^{-1}$. The stabilizer of $X_0$ is
$L=G\times_QG\cong G\times C_2$, with $(g,e)\mapsto(g,gz^e)$.
Use the coordinates
\[
 \begin{gathered}
 p=(1,0,0),\quad s=(0,0,1),\quad z=(ps)^2=(1,1,0),\\
 P=(p,0),\quad S=(s,0),\quad Z=(z,0),\quad E=(1,1)\in L.
 \end{gathered}
\]
An arbitrary pointed-bimodule cochain is $\beta'=\beta\gamma$, where
$\beta$ is a fixed cochain for the forced associator and $\gamma$ is an
arbitrary ordinary $\mathbb C^\times$-valued two-cocycle on $L$.
For commuting $A,B$ write
$C_\gamma(A,B)=\gamma(A,B)/\gamma(B,A)$.
Let $T$ exchange the two $K$ factors and let $\phi_q$ be conjugation
by $(q,1)$ on the stabilizer. Categorical duality with shift $qH$ requires
\[
 c'_q=\phi_q^*(\beta')\,\kappa_q\,T^*(\beta')
\]
to be a coboundary. Here $\kappa_q$ is the normalizer correction for the
opposite-bimodule convention of Vainerman--Vallin, Sections 2.2.1, 2.2.2
and 3.1.1. No grading of the whole rank-twelve category is assumed.
The explicit cochain and correction are recorded in
\path{evidence/unresolved10_20261027/r9_noncommutative/REPORT.md}; the short
cochain verifier checks every coboundary equation and the identities below.

For $q=(x,y,n)$ their commuting-pair identity is
\[
 C_{c'_q}(P,E)C_{c'_q}(S,E)=(-1)^{x+y+n}.
\]
It excludes both reflection cosets for every complex $\gamma$.
For either lift of the remaining rotation coset, represented by
$q=ps=(0,1,1)$, direct substitution gives
\[
 C_{c_q}(P,E)=1,\qquad C_\beta(P,Z)=-1.
\]
The change of the first expression on replacing $\beta$ by $\beta\gamma$
is exactly $C_\gamma(P,Z)$. To see this without a root-of-unity bound,
centrality of $E$ makes $C_\gamma(-,E)$ a character; it kills $Z=[P,S]$
and has square one on $P$. For some $a\in\mathbb F_2$,
\[
 \phi_q(P)=P(ZE)^a,\quad\phi_q(E)=E,\quad TP=P,\quad TE=ZE.
\]
Consequently
\[
 C_{\phi_q^*\gamma}(P,E)C_{T^*\gamma}(P,E)
 =C_\gamma(P,E)^2 C_\gamma(P,Z)=C_\gamma(P,Z).
\]
A coboundary has trivial commutator on commuting pairs. The duality
condition therefore forces $C_\gamma(P,Z)=1$ and hence
\begin{equation}\label{eq:r1975-projective-sign}
 C_{\beta'}(P,Z)=-1.
\end{equation}
This holds for arbitrary complex cochains and gauges. The rotation-duality
equation itself has an eighth-root trivialization; its consistency is not
the contradiction, and a fourth-root-only search would be invalid.

\subsubsection{Every central restriction has even multiplicity}
The chosen pointed associator is identically one on
$A=\langle p,z\rangle\cong V_4$. Diagonal conjugation by $A$ therefore
has its strict pointed convention, without a global normalization of
the diagonal $D_8$ tensorator. For any $W\in Z(\mathcal D)$,
its half-braiding gives $F(W)$ an $A$-equivariant structure. Conjugation
fixes every $X_q$, so $V=\operatorname{Hom}(X_0,F(W))$ carries the
projective stabilizer action from $\beta'|_A$, or its inverse. Either
convention has the commutator in \eqref{eq:r1975-projective-sign}.
Thus its two invertible operators obey $U_pU_z=-U_zU_p$.
Taking determinants gives $(-1)^{\dim V}=1$. We obtain
\begin{equation}\label{eq:r1975-even-restriction}
 [F(W):X_0]\ \text{is even for every }W\in Z(\mathcal D).
\end{equation}
This step uses neither indicator signs nor unitarity.

\subsubsection{The distinguished Galois object has odd multiplicity}
Pass to the pseudo-unitary Galois representative supplied by the signed
existence gate, and take the actual Morita dual. Both categories are
pseudo-unitary because their intrinsic global dimensions and FP dimensions
are Morita invariant. This does not assert pseudo-unitarity of an arbitrary
original category. Let $\sigma(\sqrt2)=-\sqrt2$ and $D'=\sigma(D)$.
The uniform ring has eight scalar irreducibles and one irreducible of
degree two. The two dimension characters have codegrees $D,D'$.
For each nontrivial character $\psi$ of $Q$, there are two representations
$g\mapsto\psi(gH)$, $X_q\mapsto t\psi(q)$, $t^2=2\psi(\alpha)$;
their codegrees are $16$. The two-dimensional representation of $D_8$
extends by $X_q\mapsto0$ and has codegree $4$. Its Casimir eigenvalue
is $2\cdot4=8$, not $4$. The degree sum $8+2^2=12$ proves completeness.

The central element $v=1+\sum_qX_q$ acts on every irreducible $E$ by
$f_E/\sigma(f_E)$: it is one on the rational-codegree representations,
and $1+4d=D/D'$ or $1+4\sigma(d)=D'/D$ on the dimension characters.
Let $A_E\subset I(1)$ be the distinguished central simple of dimension
$D/f_E$ and let $B=A_{\sigma\mathrm{FP}}$. With normalized central
character $\lambda_E=\chi_E/\dim(E)$, symmetry of the unnormalized
modular $S$ matrix gives
\[
 \frac D{f_E}\lambda_E(F(B))=S_{A_E,B}=S_{B,A_E}
 =\frac D{D'}\,\sigma(D/f_E)=\frac D{\sigma(f_E)}.
\]
The degree normalization is retained. Semisimplicity of the based algebra
therefore identifies $F(B)$ with $v$:
\[
 F(B)=1+X_0+X_1+X_2+X_3.
\]
This contradicts \eqref{eq:r1975-even-restriction}. Reflection, rotation
and the previously eliminated extreme branch exhaust the possibilities,
so $R9\_2\_1975$ has no fusion categorification over $\mathbb C$.
The distinguished-object and central-character inputs are Ostrik,
Theorem 2.13 and Section 2.3, and Deshpande, Proposition 2.8 and
Section 2.7. The general-complex signed gate and both source tensor
identities have separate exact completeness replays.

\hypertarget{frontier-two-rank-nine-exclusions-from-a-symmetric-galois-unit-pairing}{%
\subsection{Two rank-nine exclusions from a symmetric Galois-unit
pairing}\label{frontier-two-rank-nine-exclusions-from-a-symmetric-galois-unit-pairing}}

\hypertarget{frontier-result-and-scope}{%
\subsubsection{Result and scope}\label{frontier-result-and-scope}}

\textbf{R9\_2\_1607 and R9\_2\_1976 have no fusion categorification over
the complex numbers.} R9\_2\_137 remains unresolved; this argument
forces all four of its large simple objects to have second
Frobenius--Schur indicator \(+1\) in the positive representative.

The argument combines the new reflection-duality obstruction in
\path{../r9_noncommutative/REPORT.md}, the Galois-unit calculation in
\path{../r9_noncommutative/GALOIS_PARITY_EXCLUSION.md}, and the
symmetric pairing observation below. It uses the previously audited
categorical TY/Morita reductions and the
general-complex-to-pseudo-unitary existence gate. There is no unitarity
assumption and no search over associators.

\hypertarget{frontier-exact-source-bindings}{%
\subsubsection{1. Exact source bindings}\label{frontier-exact-source-bindings}}

Inputs are the tensors in \path{data/current_remaining_cases.py} at
the frozen source commit
\texttt{6a426acf6465dd674351be92dcca41b72f371521}.

\begin{center}\begin{tabular}{lcll}\toprule
Ring & TY object & Large objects & Selfdual large objects\\\midrule
$R9\_2\_137$ & 8 & 3,4,5,6 & all four\\
$R9\_2\_1607$ & 3 & 1,2,7,8 & 7,8\\
$R9\_2\_1976$ & 5 & 1,2,3,4 & none\\\bottomrule
\end{tabular}\end{center}
The full tensor hashes are recorded in the source-bound accepted ledger.


Every source has four invertibles, one TY object \(Y\) of dimension 2,
and four large objects of dimension

\[
d=4+3\sqrt2,\qquad d^2=2+8d,\qquad
D=144+96\sqrt2=8+4d^2.
\]

The short verifier \texttt{verify\_source\_checks.py} independently
checks these dimensions, dualities, and hashes. Its receipt is
\texttt{source\_checks.json}.

The prior reductions establish a connected subgroup algebra and a
rank-twelve Morita dual with invertible group \(G=D_8\), central
stabilizer \(H=C_2\), quotient \(Q=G/H=V_4\), and four other simples
\(X_q\). The earlier possibilities for these three sources had zero or
reflection duality shift. The new all-complex commuting-pair obstruction
excludes the reflection shifts. Thus a remaining categorification must
have the uniform zero-shift rules

\[
gX_q=X_qg=X_{q+gH},\qquad
X_pX_q=\sum_{gH=p+q}g+2\sum_{t\in Q}X_t,
\qquad X_q^*=X_q.
\]

The existing spherical/Galois existence gate allows us to work with a
pseudo-unitary representative of this Morita pair. This does not assert
that the original hypothetical category was unitary. Pseudo-unitarity is
preserved under categorical Morita equivalence, so both sides have their
canonical positive spherical structures.

\hypertarget{frontier-a-multiplicity-one-galois-object-is-symmetrically-selfdual}{%
\subsubsection{2. A multiplicity-one Galois object is symmetrically
selfdual}\label{frontier-a-multiplicity-one-galois-object-is-symmetrically-selfdual}}

Let \(\mathcal D\) be the hypothetical positive rank-twelve category,
\(F:Z(\mathcal D)\to\mathcal D\) the forgetful functor, and \(I\) its
right adjoint. Its canonical algebra \(L=I(\mathbf1)\) is connected
étale and has trivial twist. It has a nondegenerate Frobenius pairing. A
commutative Frobenius algebra with trivial twist in a ribbon category is
symmetric in the \textbf{pivotal} sense: see
Fröhlich--Fuchs--Runkel--Schweigert, \emph{Correspondences of ribbon
categories}, Proposition 2.25(ii), https://arxiv.org/pdf/math/0309465 or
the author copy
https://www.math.uni-hamburg.de/home/runkel/PDF/corr.pdf.

In explicit duality notation, its pairing gives an isomorphism
\(\phi_L:L\to L^*\) satisfying

\[
\phi_L=\phi_L^*j_L,
\]

where \(j\) is the pivotal structure. This is stronger than merely
braided commutativity of the multiplication; it is the property used
here.

Let \(B=A_{\sigma\mathrm{FP}}\) be the distinguished summand of \(L\)
for the conjugate dimension character, where \(\sigma(\sqrt2)=-\sqrt2\).
For the zero-shift uniform ring the complete list of irreducibles
consists of eight one-dimensional representations and one
two-dimensional representation. Their formal codegrees are

\[
D,\quad D'=\sigma(D),\quad 16\text{ (six times)},\quad4.
\]

The distinguished central objects therefore have dimensions

\[
1,\quad D/D',\quad D/16\text{ (six times)},\quad D/4.
\]

Their multiplicities in \(L\) are respectively 1 for the scalar
representations and 2 for the two-dimensional representation. In
particular \(B\) occurs once and its dimension \(D/D'\) is unique among
these constituents. Since \(L\cong L^*\), duality must preserve this
unique dimension and consequently \(B^*\cong B\).

The pairing of \(L\) cannot pair this summand with any other simple
summand. Its restriction to \(B\) is therefore nondegenerate and still
pivotally symmetric. Thus \(B\) is symmetrically selfdual.

The normalized central-character calculation, with the factor
\(\dim(E)\) retained for the two-dimensional representation, gives

\[
F(B)=\mathbf1+\sum_{q\in Q}X_q.
\tag{1}
\]

For a full proof see the sibling \texttt{GALOIS\_PARITY\_EXCLUSION.md},
§4. The same proof applies with zero shift: the rational codegrees
remain 16 and 4. More explicitly, \(1+\sum X_q\) acts on every
irreducible \(E\) by \(f_E/\sigma(f_E)\), and the modular \(S\)-symmetry
formula gives exactly the same scalar for \(F(B)\). Equality follows
from semisimplicity of the complexified based algebra.

The functor \(F\) is pivotal and strong monoidal. It therefore carries
the nondegenerate symmetric pairing of \(B\) to one on \(F(B)\). Every
\(X_q\) is selfdual and appears once in (1). Restricting the pairing to
each of these pairwise nonisomorphic simple summands is again
nondegenerate and symmetric. Hence

\[
\boxed{\nu_2^{\mathcal D}(X_q)=+1\quad\text{for every }q\in Q.}
\tag{2}
\]

We are \textbf{not} claiming that numerical Frobenius--Schur indicators
are preserved by arbitrary tensor functors. The transported object is a
particular symmetric pairing, and multiplicity one is essential.

\hypertarget{frontier-exact-morita-invariant-second-indicator-comparison}{%
\subsubsection{3. Exact Morita-invariant second-indicator
comparison}\label{frontier-exact-morita-invariant-second-indicator-comparison}}

For a pseudo-unitary fusion category define

\[
T_2(\mathcal C)=\sum_{V\in\mathrm{Irr}(\mathcal C)}
\dim(V)\nu_2(V).
\]

The center-induction indicator formula gives

\[
T_2(\mathcal C)=\frac1{\dim\mathcal C}
\sum_{Z\in\mathrm{Irr}(Z(\mathcal C))}\theta_Z^2\dim(Z)^2.
\tag{3}
\]

Indeed, multiply
\(\nu_2(V)=\operatorname{Tr}(\theta_{I(V)}^2)/\dim\mathcal C\) by
\(\dim V\), sum over \(V\), and use Frobenius reciprocity:
\(\sum_V\dim(V)[I(V):Z]=\dim F(Z)=\dim Z\). Formula (3) is invariant
under Morita equivalence because the centers are braided equivalent and
their canonical positive ribbon structures agree. The indicator formula
is Ng--Schauenburg, \emph{Frobenius--Schur indicators and exponents of
spherical categories}, https://arxiv.org/pdf/math/0601012.

Write the TY sign in the source as \(\varepsilon\in\{+1,-1\}\), so
\(\nu_2(Y)=\varepsilon\). Its four invertibles have trivial associator
and second indicator 1. Thus its integral contribution is

\[
T_{2,\mathrm{integral}}(\mathcal C)=4+2\varepsilon.
\]

In the rank-twelve dual, the forced pointed associator is trivial on the
even Klein four subgroup, while the two odd reflections have second
indicator \(\varepsilon\). The two order-four rotations are not selfdual
and have second indicator zero. Thus its pointed contribution is also
exactly

\[
T_{2,\mathrm{pointed}}(\mathcal D)=4+2\varepsilon.
\]

These signs are part of the earlier TY-origin associator calculation;
Galois conjugation does not change them because they are \(+1,-1\).
Using (2) and cancelling these equal integral contributions in (3) gives

\[
\boxed{\sum_{\text{four large }V\text{ in }\mathcal C}\nu_2(V)=4.}
\tag{4}
\]

Now a nonselfdual simple has second indicator zero, whereas a selfdual
simple has second indicator \(+1\) or \(-1\):

\begin{itemize}
\tightlist
\item
  \textbf{R9\_2\_1607:} only two large simples are selfdual. Their total
  is at most 2, contradicting (4).
\item
  \textbf{R9\_2\_1976:} none of the four large simples is selfdual.
  Their total is zero, contradicting (4).
\item
  \textbf{R9\_2\_137:} all four are selfdual. Equation (4) forces all
  four indicators to be \(+1\), but supplies no contradiction.
\end{itemize}

This completes the two exclusions over \(\mathbb C\), using the stated
general-complex existence gate. No numerical indicator optimization,
bounded root search, or unitarity classification enters the new
argument.

\subsection{The remaining ring $R9\_2\_137$}
In the remaining zero-shift Morita category, put $S=\sum_q X_q$ and let
$z$ be the central pointed involution. The pointed-subcategory canonical
central algebra is $1+B$ and its augmentation splits the forgetful algebra
as $1\times A$, with $[A]=1+S$. Thus A is forced in a hypothetical category.
Its induction Gram gives eight modules underlying $g+S$ and four underlying
$X_q$. Each of the latter has internal end $1+z$. The module category is
therefore the central-involution module already used, not a new construction.
This identifies the forced algebra without constructing a new source category.
The strongest necessary center systems, incorporating whole-column Galois
orbits and fourth rotations, are treated in Section~\ref{sec:r137-restrictions}.
Both TY signs remain feasible there; their witnesses are not categorifications.
